\section{Introduction}
\label{sec:intro}

\subsection{What the project does}

The project reads one printed line of text or mathematics with a small camera and reproduces it
as a typeset document. The camera is the OV7675 module on the USC EE\,105 board, which carries an
Arduino Nano 33 BLE Sense Lite. The board has too little memory to do more than take the
photograph (256\,KB of RAM against 76{,}800 bytes for a single $320 \times 240$ grayscale frame),
so it sends each frame to a laptop over the USB cable, and the laptop does the rest. The sequence
of processing steps on the laptop is called the \term{pipeline}; Figure~\ref{fig:pipeline} shows
it.

\begin{figure}[H]
\centering
\begin{tikzpicture}[node distance=7mm, every node/.style={font=\small},
  box/.style={draw, rounded corners=2pt, align=center, minimum height=12mm, inner sep=2.5pt, text width=23mm},
  >={Stealth[length=2mm]}]
\node[box, fill=black!6] (b) {Board\\\scriptsize camera, sketch};
\node[box, right=of b] (c) {Capture\\\scriptsize 320\,$\times$\,240 gray};
\node[box, right=of c] (s) {Segment\\\scriptsize one box per symbol};
\node[box, right=of s] (k) {Classify\\\scriptsize CNN, 159 classes};
\node[box, below=8mm of k] (l) {Layout\\\scriptsize scripts, fractions};
\node[box, left=of l] (t) {\LaTeX{} line\\\scriptsize compile check};
\node[box, left=of t] (p) {PDF\\\scriptsize pdflatex};
\draw[->] (b) -- node[above, font=\tiny, inner sep=1.5pt]{USB} (c);
\draw[->] (c) -- (s);
\draw[->] (s) -- (k);
\draw[->] (k) -- (l);
\draw[->] (l) -- (t);
\draw[->] (t) -- (p);
\end{tikzpicture}
\caption{The pipeline. Only the shaded box runs on the board. This guide installs and tests every
unshaded box, using stored photographs in place of the board.}
\label{fig:pipeline}
\end{figure}

The stages are as follows. \emph{Capture} receives a frame from the board. \emph{Segment} finds
the ink in the frame and draws one box around each printed symbol. \emph{Classify} passes each
boxed symbol to a convolutional neural network (CNN), which names it as one of 159 classes
covering the visible ASCII characters and the common mathematical symbols. \emph{Layout} decides
how the symbols relate to one another (which are superscripts, subscripts, numerators or
denominators) and writes the corresponding \LaTeX{} text. Finally the program \cmd{pdflatex}
turns that text into a PDF.

\subsection{What you will have at the end}

A reader who completes every checkpoint of this guide has:
\begin{itemize}
\item a GitHub account that can read from, and write to, the team's repository;
\item the tools Git, Python 3.12, MiKTeX and Poppler installed and verified;
\item the project in the folder \file{C:\bs ee105\bs camera-to-latex}, with its own Python
  environment;
\item a passing run of the project's automatic tests;
\item a trained classifier and a measured accuracy on 80 simulated camera photographs;
\item a PDF, opened on screen, that the pipeline produced from a photograph of an equation.
\end{itemize}
At that point the laptop is ready for the tasks that use the board.

\subsection{Assumptions and scope}

The guide makes four assumptions and no others.
\begin{enumerate}
\item The laptop runs 64-bit Windows 10 or Windows 11 on an Intel or AMD processor.
  Section~\ref{sec:prelim} shows how to check this.
\item You can install software on it, which means that you know the password Windows asks for
  when a program requests permission to make changes.
\item It has an internet connection and at least 10\,GB of free disk space.
\item None of the project's software is installed, and you have not used any of it before.
\end{enumerate}
A team member whose laptop runs macOS or Linux should speak to the project lead before starting:
the project works on those systems (the short instructions are in the file \file{README.md}), but
the installation steps in Sections~\ref{sec:powershell} to~\ref{sec:python} are written for
Windows only.

If part of the software is already installed, do not skip the corresponding step; carry out its
verification instead (each installation step ends with one). If the verification succeeds, move
on.

\begin{note}[Status of this revision]
The commands of the project itself (Sections~\ref{sec:python} and~\ref{sec:quick}) were run by
the project lead on a Windows laptop on 23 November 2026, and their outputs are quoted from that
run. The installation steps of Sections~\ref{sec:tools} and~\ref{sec:clone} follow the
publishers' instructions (listed in the Sources appendix) and the lead's experience, but this
revision has not yet been followed from its first step to its last on a second, empty laptop.
The first team member to do so is asked to report every point at which the screen differed from
the description, so that the guide can be corrected.
\end{note}

\subsection{How to read this guide}
\label{sec:conventions}

The guide is meant to be followed in order, one step at a time, with this document open beside
the windows in which you work. Steps are numbered within their section: Step~5.3 is the third
step of Section~5. Four kinds of box appear.

\begin{ps}
Get-Date
\end{ps}

A blue box holds \emph{exactly one command} for Windows PowerShell, the program introduced in
Section~\ref{sec:powershell}. Type the line exactly as printed, or copy it from this PDF and
paste it, and then press \key{Enter}. Upper and lower case, spaces, hyphens, quotation marks and
backslashes all matter. A command is always printed on a single line; if two blue boxes follow
each other, they are two commands, and the first must finish before the second is typed.

\begin{out}
Tuesday, November 24, 2026 9:41:07 AM
\end{out}

A gray box shows what appears on the screen after the command. Your screen will rarely match it
character for character (dates, version numbers and folder names differ); the text below each
gray box says which part matters.

\begin{iffail}
An amber box describes the most likely way in which a step goes wrong, and the remedy.
Section~\ref{sec:trouble} collects all of them, with the exact error messages, in one place.
\end{iffail}

\begin{note}
A green box gives background that is useful but not needed to carry out the step.
\end{note}

Three further conventions are used throughout.
\begin{itemize}
\item \textbf{Keys} are drawn like this: \key{Enter}, \key{Ctrl}+\key{C} (hold the first key and
  press the second).
\item \textbf{Things to click} are printed in bold sans-serif type, for example \ui{Sign up}, and
  a path through menus is written with arrows: \menu{Settings}\then\menu{System}\then\menu{About}.
\item \textbf{Placeholders}, which you must replace with your own information, are printed in
  capitals on a yellow background, for example \ph{YOUR-NAME}. There are only four in the whole
  guide. Nothing else in a command is ever to be replaced, and no command in this guide contains
  the characters \code|<| or \code|>| around a word.
\end{itemize}

Each step ends with a \ok~\textbf{Checkpoint}: an observable fact that tells you the step
succeeded. Do not go on to the next step until the checkpoint is true.

\subsection{When you are stuck}
\label{sec:help}

First read the amber box of the step, then the entry for your error message in
Section~\ref{sec:trouble}. If neither resolves the problem, send the project lead three things:
the step number, a screenshot of the whole PowerShell window (press
\key{Windows}+\key{Shift}+\key{S}, drag a rectangle over the window, and paste the picture into
a message with \key{Ctrl}+\key{V}), and one sentence saying what you expected. Do not retype an
error message by hand, and do not experiment with commands found elsewhere; an installation that
has been altered by trial and error is much harder to repair than one that stopped at a known
step.
